\chapter{Fine-tuning: the curriculum, the BatchNorm see-saw and mixed batches}\label{ch:finetune}
\begingroup\setlength{\emergencystretch}{5em}% long provenance paths: allow loose lines rather than overfull ones
\def\tbrk{\hfil\penalty50\hfilneg}% local: a line may break before a provenance tag, leaving the line ragged

\begin{onething}
Fine-tuning runs as stages (probe the frozen encoder, adapt it to three sensors, then train reconstruction and segmentation
together), each opened only by a gate written down in advance. On the way, the model's single BatchNorm turned out to keep a
running mean that is an exponential average dominated by the last couple of dozen batches, so sensor-pure batches made evaluation
depend on which sensor came last. Mixing every sensor into every batch, one stem pass per band list and one encoder pass for all,
fixed it.
\end{onething}

\section{Why a curriculum}\label{sec:finetune:why}

F001 (\cref{ch:pretrain}) has only seen SpectralEarth. The product reads PRISMA, EnMAP and AVIRIS-NG at their own bands, and now has
a segmentation head (\cref{ch:seghead}) that has never been trained. The obvious plan is to switch everything on at once: load F001,
add the head, feed all three sensors, train every parameter on reconstruction plus segmentation. It fails in three ways.
\begin{itemize}
  \item \textbf{You cannot tell what helped.} If detection moves, was it the new sensors, the new loss, or the new head pulling the encoder?
  \item \textbf{A fresh head wrecks a trained encoder.} A randomly initialised head produces large, meaningless gradients in its first steps,
    and they flow back into an encoder that took a hundred epochs to train.
  \item \textbf{The new skill can cost the old one.} The product's job is detection, which reads the \emph{reconstruction}. Training the shared
    encoder hard on segmentation can make it a worse reconstructor.
\end{itemize}
So the plan (\file{docs/proposals/fine-tuning-and-segmentation.md}, agreed 2026-10-16) is a \emph{curriculum}\index{curriculum}: a sequence of
stages, each changing one thing, each judged before the next starts. The code is \file{app/foundation_models/trainers/finetune_trainer.py}
with its configuration \code{SpectralStemFinetuneConfig}; runs live in \file{research/experiments/indradhanu-finetune/}, one folder per run,
each with its pass criteria written before launch. This chapter is about the machinery and its mathematics; what each run showed is told
in \cref{ch:ftruns}.

\section{The stages and their gates}\label{sec:finetune:stages}

\begin{table}[htbp]
\centering\small
\caption{The fine-tuning curriculum. Stage 4 is optional.}
\label{tab:finetune:stages}
\begin{tabular}{p{1.9cm}p{1.1cm}p{2.9cm}p{2.6cm}p{3.9cm}}
\toprule
stage & mode & trains & data & gate before the next\\
\midrule
0 setup & -- & nothing & -- & everything exists and is tested; label policy fixed\\
1 probe (FT01) & \texttt{probe} & the segmentation head only & EnMAP & per-class IoU recorded\\
2 adapt (FT02) & \texttt{adapt} & stem, encoder, reconstruction decoder & all three, by share & bench $\ge 0.861$; test reconstruction falls on every sensor\\
3 joint (FT03) & \texttt{joint} & everything & all three; segmentation loss on EnMAP & bench not below stage 2; IoU above stage 1\\
4 pseudo-labels & \texttt{joint} & as stage 3 & plus confident PRISMA and AVIRIS-NG labels & IoU holds; bench holds\\
5 product & -- & -- & -- & masks replace the threshold rules\\
\bottomrule
\end{tabular}
\end{table}

\begin{definition}[Bench and gate]\index{gate (fine-tuning)}
The \emph{bench} is the mean ROC-AUC (\cref{ch:measures}) on the four labelled PRISMA scenes under the default score, residual Mahalanobis
(\cref{ch:residualmaha}). A \emph{gate} is a pass criterion written in a run's README before launch and never moved afterwards. 0.861 is F001 at
epoch 60.
\end{definition}

\coursefig[0.95\linewidth]{iw-49-curriculum}{(a) The six stages, what each trains, and the gate (amber) it must pass; stage 4 is optional. (b) Batches
are drawn by fixed shares (top), not in proportion to each sensor's tile count (bottom), which would give AVIRIS-NG 43\% of training.}{fig:finetune:curriculum}

\paragraph{Why the gate is read at \texttt{best.pt}, chosen by another rule.} A gate is read at the checkpoint that validation reconstruction picks
(\cref{ch:pretrain}), never at whichever epoch scores best on the bench. Suppose a run's true detection ability does not change at all, and four
checkpoints read 0.855, 0.858, 0.864 and 0.853 on the four bench scenes, the differences being noise\tbrk\tagA. Picking the best of the four reports 0.864,
which passes a 0.861 gate, although nothing improved: the maximum of several noisy readings is biased upwards. Choosing by the bench turns the bench
into training data. \Cref{ch:bootstrap} measures how large that noise is.

\section{Starting from F001}\label{sec:finetune:init}

\paragraph{Loading with a check.} \code{build_model} loads the pretrained checkpoint into a model that now has a head:
\begin{lstlisting}
missing, unexpected = model.load_state_dict(state, strict=False)
# new_heads = ("seg_head.", "scale_head."): the only tensors allowed to be missing
if unexpected or any(not k.startswith(new_heads) for k in missing):
    raise ValueError(f"init_from does not match the model: ...")
\end{lstlisting}
Everything except the new heads must match exactly. A configuration that builds the wrong architecture fails here, loudly, instead of training half
a model from scratch.

\paragraph{The normalisation curve comes from the checkpoint.} Inside the model each band is normalised as $(v - \mu(\lambda))/\sigma(\lambda)$ with the
curve computed on SpectralEarth (\cref{ch:pretrain}). The fine-tuning configuration takes that curve from \code{init_from}'s checkpoint and never
recomputes it.

\begin{worked}[label={we:finetune:curve}]{Why the curve is frozen}
One band with pretraining mean 0.20 and spread 0.10, and a pixel reading 0.25. Suppose the curve were recomputed on EnMAP and came out at mean 0.22,
spread 0.12 (illustrative\tbrk\tagA).
\begin{align*}
  \text{pretraining curve:}&\quad (0.25 - 0.20)/0.10 = 0.05 / 0.10 = 0.50\\
  \text{recomputed curve:}&\quad (0.25 - 0.22)/0.12 = 0.03 / 0.12 = 0.25.
\end{align*}
The same pixel would reach the weights as 0.25 instead of the 0.50 they learned to read, and every input in every band would shift at once.
\end{worked}

\section{Parameters, buffers and BatchNorm}\label{sec:finetune:bn}

\subsection{Freezing parameters}

A \emph{probe}\index{probe} trains a small head on top of frozen features, to measure what the features already know. If water and cloud separate well
with the encoder untouched, F001 learned them in pretraining. In probe mode the trainer marks every parameter outside the head as not trainable and
gives the optimiser only the head's parameters:
\begin{lstlisting}
for name, p in self.model.named_parameters():
    p.requires_grad = name.startswith("seg_head.")       # only the head learns
self.optimizer = torch.optim.Adam([p for p in self.model.parameters() if p.requires_grad],
                                  lr=config.learning_rate)
\end{lstlisting}
(Adam is the adaptive optimiser of \cref{ch:trainloop} \citep{kingma2015adam}, here without decoupled weight decay.) That is not quite enough, because
a model holds two kinds of tensors.\index{parameter (model)}
\begin{definition}[Parameter and buffer]\index{buffer (PyTorch)}
A \emph{parameter} is a tensor the optimiser changes through gradients. A \emph{buffer} is a tensor the model keeps and may change itself, without
gradients; \code{requires_grad} does not touch it.
\end{definition}

\subsection{BatchNorm: two modes}

The model has one \emph{batch normalisation}\index{batch normalisation} layer \citep{ioffe2015batch}, \code{self.norm = nn.BatchNorm2d(D)} with
$D = 32$, between the spectral stem and the encoder (\file{app/foundation_models/components/spectral_stem_seg_former_mae.py}, line 75). Every pixel's 32
stem channels pass through it. The encoder's blocks use LayerNorm (\cref{ch:block}), which normalises each token on its own and keeps no statistics.
So this is the model's only layer whose output depends on the \emph{other} samples in a batch.

For one channel, with $x$ the values of that channel over every pixel of every sample in the batch:
\begin{description}[style=nextline,leftmargin=1.2em]
  \item[Training mode (\texttt{model.train()})] Compute the batch's mean $\mu_B$ and variance $\sigma_B^2$ over all those values, and output
    \begin{equation}\label{eq:finetune:bntrain}
      y = \gamma\,\frac{x - \mu_B}{\sqrt{\sigma_B^2 + \epsilon}} + \beta,
    \end{equation}
    with learned scale $\gamma$ and shift $\beta$ (parameters) and $\epsilon = 10^{-5}$. Also update two buffers, the running mean $r$ and running
    variance, towards the batch's values:
    \begin{equation}\label{eq:finetune:ema}
      r \leftarrow (1 - m)\,r + m\,\mu_B, \qquad m = 0.1.
    \end{equation}
  \item[Evaluation mode (\texttt{model.eval()})] Use the stored running mean and variance in place of $\mu_B$ and $\sigma_B^2$, and update nothing.
\end{description}
\code{FoundationTrainer._run_epoch} (\file{app/abstract_classes/foundation_trainer.py}, lines 213 and 228) calls \code{model.train()} before the
training pass and \code{model.eval()} before validation; inference also runs in evaluation mode. Training and evaluation therefore normalise by
\emph{different numbers}. That is harmless when every batch looks like the whole data set.

\begin{sourcenote}[: the momentum]
\code{nn.BatchNorm2d(D)} is built with no \code{momentum} argument (line 75), so PyTorch's default $m = 0.1$ applies, with the convention of
\eqref{eq:finetune:ema}: the \emph{new} batch gets weight $m$. (Some other libraries call $1 - m = 0.9$ the momentum; read the formula, not the name.)
\file{docs/lld/finetune.md} records the same: ``PyTorch default momentum 0.1''. One detail the courses leave out: PyTorch updates the running
\emph{variance} with the batch's unbiased variance (divided by $n - 1$) while normalising the batch itself with the biased one (divided by $n$); with
$8 \times 128 \times 128 = 131{,}072$ values per channel per micro-batch, the difference is a factor $131{,}072/131{,}071$, negligible.
\end{sourcenote}

\subsection{The running mean is an exponential moving average}

Unroll \eqref{eq:finetune:ema}. Write $x_k$ for the $k$-th batch's mean and $r_0$ for the starting value; $0.9 = 1 - m$.
\begin{align*}
  r_1 &= 0.9\,r_0 + 0.1\,x_1\\
  r_2 &= 0.9\,r_1 + 0.1\,x_2 = 0.9(0.9\,r_0 + 0.1\,x_1) + 0.1\,x_2 = 0.9^2 r_0 + 0.1(0.9\,x_1 + x_2)\\
  r_3 &= 0.9\,r_2 + 0.1\,x_3 = 0.9^3 r_0 + 0.1(0.9^2 x_1 + 0.9\,x_2 + x_3).
\intertext{The pattern: each older batch's weight is multiplied by another 0.9. After $n$ batches}
  r_n &= 0.9^{n}\,r_0 + 0.1\sum_{k=1}^{n} 0.9^{\,n-k}\,x_k.
\end{align*}
This is an \emph{exponential moving average}\index{exponential moving average}: the newest batch has weight 0.1, the one before 0.09, then 0.081, and so on.

\paragraph{How much weight do the last $K$ batches carry?} Add the weights of batches $n, n-1, \dots, n-K+1$:
\begin{align*}
  W_K &= 0.1\,(1 + 0.9 + 0.9^2 + \dots + 0.9^{K-1}).
\intertext{The bracket is a \emph{geometric series}\index{geometric series}. Check the formula $1 + q + \dots + q^{K-1} = (1 - q^K)/(1 - q)$ on a tiny case, $q = 0.9$, $K = 3$:
$1 + 0.9 + 0.81 = 2.71$ and $(1 - 0.729)/(1 - 0.9) = 0.271/0.1 = 2.71$. So}
  W_K &= 0.1 \cdot \frac{1 - 0.9^{K}}{1 - 0.9} = 0.1 \cdot \frac{1 - 0.9^K}{0.1} = 1 - 0.9^{K}.
\end{align*}
\begin{center}
\begin{tabular}{rll}
\toprule
$K$ & $0.9^K$ & weight of the last $K$ batches, $1 - 0.9^K$\\
\midrule
1 & 0.9 & 0.10\\
5 & 0.590 & 0.41\\
10 & 0.349 & 0.651\\
22 & 0.098 & 0.902\\
\bottomrule
\end{tabular}
\end{center}
How many batches make 90\%? Solve $1 - 0.9^K = 0.9$, that is $0.9^K = 0.1$. Take logarithms: $K\ln 0.9 = \ln 0.1$, so
$K = \ln 0.1 / \ln 0.9 = -2.3026 / -0.1054 = 21.85$. About 22 batches carry 90\% of the stored mean.

\begin{worked}[label={we:finetune:drift}]{A ``frozen'' encoder that drifts}
After pretraining a channel's running mean is 0.00, and EnMAP batches have mean 0.50 (illustrative\tbrk\tagA). The probe freezes every parameter but leaves
the model in training mode. With all $x_k = 0.5$ and $r_0 = 0$, the unrolled formula gives $r_n = 0.5\,(1 - 0.9^n)$:
\begin{align*}
  \text{after 1 batch:}&\quad 0.9 \times 0.00 + 0.1 \times 0.50 = 0.050\\
  \text{after 2 batches:}&\quad 0.9 \times 0.050 + 0.1 \times 0.50 = 0.045 + 0.050 = 0.095\\
  \text{after 10 batches:}&\quad 0.50 \times (1 - 0.349) = 0.50 \times 0.651 = 0.326.
\end{align*}
Ten batches in, the ``frozen'' model normalises its input differently from the one that was pretrained, and the probe measures a drifted encoder.
\end{worked}

So the probe puts the whole model in evaluation mode for the training pass and switches only the head back:
\begin{lstlisting}
def _train_pass_mode(self) -> None:
    mode = self.config.model_config_.mode
    # frozen backbone: keep its BatchNorm statistics
    if mode in ("probe", "scale", "decoder"):
        self.model.eval()
        {"probe": self.model.seg_head, ...}[mode].train()
\end{lstlisting}
It runs after the parent's \code{model.train()} and before the training pass. Evaluation mode also turns off dropout (\code{drop_rate} 0.4 in F001), so
the probe sees exactly the features inference sees. In probe mode the validation score that ranks epochs is the segmentation loss, since nothing else
is trained, and a batch without labels contributes nothing, which is why the probe reads EnMAP alone.

\section{Mixing sensors by share}\label{sec:finetune:shares}

FT01 and FT02 drew every batch from one sensor (the stem took one band list per call, \cref{ch:pretrain}). \code{_MixedLoader} chooses the sensor for
each batch at random, weighted by its share:
\begin{lstlisting}
name = self.rng.choices(self.names, self.weights)[0]     # weights e.g. 0.40, 0.35, 0.25
\end{lstlisting}
It keeps one iterator per sensor, restarts a sensor whose shards run out, and seeds its generator with the starting epoch, so a resumed run draws a
fresh but reproducible sequence.

\paragraph{Why shares, not one shuffled pool.} Count the training tiles:\tbrk\tagM{research/experiments/indradhanu-finetune/README.md}
\begin{align*}
  \text{total} &= 14{,}542 + 14{,}802 + 22{,}377 = 51{,}721\\
  \text{EnMAP} &= 14{,}542 / 51{,}721 = 0.281\\
  \text{PRISMA} &= 14{,}802 / 51{,}721 = 0.286\\
  \text{AVIRIS-NG} &= 22{,}377 / 51{,}721 = 0.433.
\end{align*}
A pool would give AVIRIS-NG 43\% of training: one airborne campaign, about 200 tiles per flight line, against the product's two satellites. The
starting shares are EnMAP 0.40, PRISMA 0.35, AVIRIS-NG 0.25.

\paragraph{How exact are random shares?} Each batch is EnMAP with probability 0.40, independently, so over $n$ batches the EnMAP count is a
\emph{binomial}\index{binomial distribution} variable with mean $np$ and standard deviation $\sqrt{np(1-p)}$. Over 1,000 batches:
\begin{align*}
  \text{mean} &= 1{,}000 \times 0.40 = 400\\
  \text{standard deviation} &= \sqrt{1{,}000 \times 0.40 \times 0.60} = \sqrt{240} = 15.5.
\end{align*}
So about 400 $\pm$ 15 EnMAP batches per thousand: shares are honoured to a few percent over an epoch of 3,750 micro-batches.

\paragraph{A validation trap.} Validation runs every test source to the end, in order (\code{_SequentialLoader}), and
\code{test_samples_per_epoch} caps the total. The test sets hold $950 + 1{,}570 + 2{,}270 = 4{,}790$ tiles. A cap below 4,790 cuts off whichever sensor
comes last, so its per-sensor score would rest on part of its test set or none of it. FT02 and FT03 set the cap to 5,000.\tbrk\tagM{research/experiments/indradhanu-finetune/FT03-joint/config.json}

\section{The trained-wavelength record}\label{sec:finetune:record}

\Cref{ch:pretrain}'s third band fix drops, at load time, any band more than 15\,nm from every wavelength the model has trained on. For F001 that list is
SpectralEarth's 202. Fine-tuning trains on PRISMA's own 234 bands, so after it the 20 PRISMA bench bands that F001 dropped are no longer untrained. The
model keeps a buffer\index{trained-wavelength record}, \code{trained_wl}, that starts as the pretraining wavelengths and grows: each time a batch arrives with a band list the trainer has
not seen, \code{extend_trained_wavelengths} adds every wavelength not within 0.5\,nm of one already recorded, counting near-duplicates among the new ones
once.

\begin{worked}[label={we:finetune:record}]{Growing the record}
Record 500, 510, 520\,nm; a new list 500.3, 515.0, 515.2, 530.0.
\begin{center}
\begin{tabular}{llll}
\toprule
new & nearest recorded & distance & result\\
\midrule
500.3 & 500 & 0.3 $\le 0.5$ & already known, skipped\\
515.0 & 510 or 520 & 5.0 $> 0.5$ & new\\
515.2 & 515.0 (new) & 0.2 $\le 0.5$ & counted once with 515.0\\
530.0 & 520 & 10.0 $> 0.5$ & new\\
\bottomrule
\end{tabular}
\end{center}
The record becomes 500, 510, 515, 520, 530, and the call returns 2. \code{trained_band_mask} then accepts a band within 15\,nm of anything recorded: 544\,nm
is $544 - 530 = 14$\,nm away and accepted; 546\,nm is 16\,nm away and dropped.
\end{worked}
The two tolerances do different jobs: 0.5\,nm decides ``is this the same wavelength?''; 15\,nm decides ``is this close enough to trained wavelengths for the
wavelength networks to read it?''. The record grows when a band list is \emph{seen}, not when it is \emph{learned}; ``trained on'' is the most the code can
know.

\begin{drift}
During FT02 the record grew from 202 to 733 wavelengths: AVIRIS-NG added 337, PRISMA 177, EnMAP 17 ($202 + 337 + 177 + 17 = 733$).\tbrk\tagM{research/experiments/indradhanu-finetune/FT02-adapt/README.md}
On the bench's 185 PRISMA bands, 165 were accepted before fine-tuning and all 185 after.
\end{drift}

\section{The see-saw and the one BatchNorm}\label{sec:finetune:seesaw}

\subsection{The symptom}

FT02 adapted F001 on all three sensors with one sensor per micro-batch (micro-batches of 8, gradients accumulated over 4). Its per-sensor validation did
not fall smoothly: in epochs 11--12 and 15--17 PRISMA improved while EnMAP and AVIRIS-NG got worse, and then the reverse.

\begin{drift}
FT02 (2026-10-16), validation SAM: at epoch 8 (\code{best.pt}) PRISMA 5.93$^\circ$, AVIRIS-NG 6.07$^\circ$; at epoch 17 PRISMA 5.04$^\circ$ (its best) and
AVIRIS-NG 8.24$^\circ$ (worse than at epoch 1).\tbrk\tagM{research/experiments/indradhanu-finetune/FT02-adapt/README.md} The gate was missed: bench 0.855
against 0.861 (\cref{ch:ftruns}).
\end{drift}

Validation tiles are fixed and read in the same order every epoch, so the data did not change between epochs: the model did. Two explanations fit a
see-saw between sensors: \emph{gradient competition} (each optimiser step leans towards whichever sensors were in it) and \emph{normalisation drift}
(something that is not a weight changes with the most recent sensor). This section is about the second, because it was measured.

\subsection{Training: each sensor normalised to itself}

Give one channel three typical values after the stem (illustrative\tbrk\tagA): EnMAP around 0.2, PRISMA around 0.6, AVIRIS-NG around $-0.4$, each with spread 1.
In a PRISMA-only batch, $\mu_B \approx 0.6$, so a typical PRISMA pixel becomes $(0.6 - 0.6)/1 = 0.0$. In an AVIRIS-NG-only batch, a typical AVIRIS-NG pixel
becomes $(-0.4 - (-0.4))/1 = 0.0$. \textbf{In training the encoder never sees the sensors' offsets}: BatchNorm removes them batch by batch.

\subsection{Evaluation: one blended estimate}

At evaluation the same PRISMA pixel is normalised by the running mean. If that is 0.2, it becomes $(0.6 - 0.2)/1 = 0.4$: an offset the encoder never saw in
training. And the running mean is not a fixed 0.2. From \cref{sec:finetune:bn}, 90\% of it comes from the last 22 micro-batches. An FT02 epoch is 30,000
samples in micro-batches of 8, so $30{,}000 / 8 = 3{,}750$ micro-batches, and the last 22 are $22 / 3{,}750 = 0.6\%$ of the epoch. Which sensors those were is
chance.

\begin{worked}[label={we:finetune:lastbatches}]{The end of an epoch moves the stored mean}
The running mean is 0.20 and the last four micro-batches are PRISMA, PRISMA, PRISMA, AVIRIS-NG. Apply $r \leftarrow 0.9\,r + 0.1\,\mu_B$ four times:
\begin{align*}
  \text{PRISMA:}&\quad 0.9 \times 0.20 + 0.1 \times 0.6 = 0.18 + 0.06 = 0.24\\
  \text{PRISMA:}&\quad 0.9 \times 0.24 + 0.1 \times 0.6 = 0.216 + 0.06 = 0.276\\
  \text{PRISMA:}&\quad 0.9 \times 0.276 + 0.1 \times 0.6 = 0.2484 + 0.06 = 0.3084\\
  \text{AVIRIS-NG:}&\quad 0.9 \times 0.3084 + 0.1 \times (-0.4) = 0.27756 - 0.04 = 0.23756.
\end{align*}
Four batches moved it from 0.20 to 0.31 and back to 0.24. Validation runs with whatever the last batch left; the next epoch ends on a different sequence.
\end{worked}

\coursefig[0.95\linewidth]{iw-51-batchnorm}{(a) The last micro-batches of an epoch, one sensor each, and the running mean they drag around; validation uses
wherever it ends. (b) In training each sensor is normalised by its own batch mean (top); at evaluation every sensor is normalised by the one running mean
(bottom), so each arrives with an offset the encoder never saw.}{fig:finetune:batchnorm}

\begin{analogy}
A teacher marks three classes' exams and rescales each to a class average of 60\% before grading, so the marks look comparable. Then the head of school
rescales everything once more, using the average of whichever class happened to be marked last. Every class's marks shift, by an amount nobody chose,
and the shift changes each term with the marking order. Marking the three classes' papers shuffled together fixes it: one average, used every time. The
per-class rescaling is training-mode BatchNorm on sensor-pure batches; the head of school is the running mean; the shuffled pile is a mixed batch.
\end{analogy}

\paragraph{Gradient accumulation does not help.}\index{gradient accumulation} FT02 accumulated gradients over 4 micro-batches before each optimiser step. That mixes sensors in the
\emph{gradient}: a step is single-sensor only if all four micro-batches came from the same sensor. The draws are independent, so the probabilities multiply:
\begin{align*}
  P(\text{all EnMAP}) &= 0.40^4 = 0.40 \times 0.40 \times 0.40 \times 0.40 = 0.0256\\
  P(\text{all PRISMA}) &= 0.35^4 = 0.1225 \times 0.1225 = 0.0150\\
  P(\text{all AVIRIS-NG}) &= 0.25^4 = 0.0625 \times 0.0625 = 0.0039\\
  P(\text{single-sensor step}) &= 0.0256 + 0.0150 + 0.0039 = 0.0445,
\end{align*}
about 4.5\% of steps. But BatchNorm's statistics are computed per micro-batch, in the forward pass, before any accumulation. Accumulation cannot mix what
BatchNorm sees.

\subsection{The per-scene test}

The FT02 README records a direct test on its epoch-20 checkpoint. The idea, known as \emph{AdaBN}\index{AdaBN} (adaptive batch normalisation)
\citep{li2016revisiting}, is to replace the stored running statistics by statistics measured on the scene being scored, then score again.

\begin{drift}
FT02 epoch 20, 2026-10-16:\tbrk\tagM{research/experiments/indradhanu-finetune/FT02-adapt/README.md} re-estimating the statistics on each PRISMA bench scene moved
them 0.35--0.76 standard deviations. Residual Mahalanobis barely moved (mean 0.853 $\to$ 0.858; objects caught 410 $\to$ 402); \code{combined} swung hard
(set1 0.943 $\to$ 0.755, set4 0.780 $\to$ 0.642).
\end{drift}

Two readings. The stored statistics are far from any one PRISMA scene's, by up to three quarters of a standard deviation, the size of the offsets in the
worked example. And the raw-error score depends on them heavily while residual Mahalanobis barely does: a BatchNorm offset shifts the reconstruction error
of a whole scene systematically, and residual Mahalanobis subtracts each pixel's local mean residual and whitens by the scene's covariance, which removes most
of a systematic shift. So BatchNorm is consistent with the see-saw but does not explain the gate miss. Per-scene re-estimation is not a product option:
\code{combined} fell, and the product scores one scene at a time.

\paragraph{The options.} The FT02 README listed two: keep BatchNorm in evaluation mode during adapt, as the probe does (the statistics stay F001's, from
SpectralEarth, matching none of the three sensors), or re-estimate them on a balanced mix before each validation (fixing what evaluation sees but not the
training-time habit of normalising each sensor to itself). The user chose a third: \textbf{mixed-sensor batches}.

\section{Mixed batches: three band lists, one encoder pass}\label{sec:finetune:mixed}

\subsection{The rule that had to bend}\index{mixed batch}

The stem reads a pixel's bands as tokens carrying their wavelengths, and its \code{forward} takes one \shape{(N,)} wavelength list. EnMAP has 224 bands,
PRISMA 234, AVIRIS-NG 372: a \shape{(3, 224, H, W)} tensor and a \shape{(3, 234, H, W)} tensor cannot be stacked. But the stem's output is
\shape{(B, 32, H, W)} whatever went in. From there on, nothing knows which sensor a pixel came from. So the rule only ever needed to be
\begin{keypoint}
one band list per \emph{stem pass}, not per batch (the invariant of \file{docs/lld/finetune.md}).
\end{keypoint}

\subsection{\texttt{forward\_mixed}}

In \file{app/foundation_models/components/spectral_stem_seg_former_mae.py}:
\begin{lstlisting}
def forward_mixed(self, groups: list, with_scale: bool = False):
    # groups: [(x, wavelengths_nm, band_valid, keep_mask, out_wavelengths_nm), ...]
    if self.stem not in WAVELENGTH_STEMS:
        raise ValueError(f"{self.stem} stem reads one band list only")
    # one stem pass per sensor, outputs stacked: (8, 32, H, W)
    z = torch.cat([self._stem(x, wl, bv) for x, wl, bv, _, _ in groups])
    # the MAE keep masks, stacked in the same order
    keep = None if groups[0][3] is None else torch.cat([g[3] for g in groups])
    z = self.norm(z)                              # ONE BatchNorm over all sensors
    features = self.encoder(z, keep_mask=keep)    # ONE encoder pass
    z_hat = self._decode(z, features, keep)       # (8, 32, H, W)
    out, start = [], 0
    for x, wl, _, _, out_wl in groups:            # split by rows, decode at own bands
        wl_out = wl if out_wl is None else out_wl
        part = self.spectral_out(z_hat[start:start + len(x)], wl_out)   # (B_g, M_g, H, W)
        out.append(self.stats.denormalize(part, wl_out))
        start += len(x)
    return out
\end{lstlisting}

\begin{sourcenote}[: the listing is the current source]
The course (IW52) shows the first version, \code{z_hat = self.decoder(self.encoder(self.norm(z), keep_mask=keep))}. The source has since gained
\code{_decode} (for the progressive decoder of FT09, which also reads the normalised stem output) and an optional \code{with_scale} output (the uncertainty head
of FT06); the listing above follows the source (lines 159--180), with the \code{with_scale} return omitted. For the SegFormer decoder, \code{_decode} is
exactly \code{self.decoder(features)}, so the two versions compute the same thing.
\end{sourcenote}

\begin{table}[htbp]
\centering\small
\caption{One mixed batch of 8 tiles of 128 $\times$ 128 through \texttt{forward\_mixed}, without band augmentation.}
\label{tab:finetune:trace}
\begin{tabular}{llll}
\toprule
step & EnMAP & PRISMA & AVIRIS-NG\\
\midrule
input & \shape{(3, 224, 128, 128)} & \shape{(3, 234, 128, 128)} & \shape{(2, 372, 128, 128)}\\
\texttt{\_stem} & \shape{(3, 32, 128, 128)} & \shape{(3, 32, 128, 128)} & \shape{(2, 32, 128, 128)}\\
\texttt{torch.cat} & \multicolumn{3}{c}{\shape{(8, 32, 128, 128)}}\\
norm $\to$ encoder $\to$ decoder & \multicolumn{3}{c}{\shape{(8, 32, 128, 128)}}\\
slice rows & 0--2 & 3--5 & 6--7\\
\texttt{spectral\_out} & \shape{(3, 224, 128, 128)} & \shape{(3, 234, 128, 128)} & \shape{(2, 372, 128, 128)}\\
\bottomrule
\end{tabular}
\end{table}

\coursefig[0.95\linewidth]{iw-52-forward-mixed}{Three sensors, three stem passes (teal), one concatenated batch through the single BatchNorm, encoder and decoder
(indigo), split by row ranges and decoded at each sensor's own wavelengths.}{fig:finetune:forwardmixed}

The fixed-band stems (\code{linear}, \code{gate}) are built for one band count, so \code{forward_mixed} refuses them. The keep masks are concatenated in the
same order as the stem outputs, so row $i$ of the batch is masked with sample $i$'s own mask.

\subsection{Where the batch comes from: rounded sub-batch sizes}

With \code{mixed_batches: true}, each sensor gets its own data loader with batch size $\max(1, \operatorname{round}(s \times 8))$ for share $s$, and
\code{_AllSourcesLoader} yields one sub-batch from every sensor per step:
\begin{align*}
  \text{EnMAP:}&\quad \operatorname{round}(0.40 \times 8) = \operatorname{round}(3.2) = 3\\
  \text{PRISMA:}&\quad \operatorname{round}(0.35 \times 8) = \operatorname{round}(2.8) = 3\\
  \text{AVIRIS-NG:}&\quad \operatorname{round}(0.25 \times 8) = \operatorname{round}(2.0) = 2\\
  \text{total:}&\quad 3 + 3 + 2 = 8.
\end{align*}
The effective shares are $3/8, 3/8, 2/8 = 0.375, 0.375, 0.25$, not 0.40, 0.35, 0.25: PRISMA gains 2.5 points and EnMAP loses 2.5. Nothing forces the rounded
sizes to sum to the batch size for other shares; at 8 with these shares they happen to.

\subsection{The loss: every tile counts equally}

\code{_mixed_loss} builds, per sensor, that sensor's MAE mask and its own band-augmentation draw, makes one \code{forward_mixed} call, and combines the
per-sensor losses weighted by kept samples $n_g$:
\begin{lstlisting}
total = total + (l1 + self._get_sam_weight(self._current_epoch) * sam + edge) * n
kept = kept + n                 # (one line in the source)
...
loss = total / kept
\end{lstlisting}
Why weight by $n_g$? If $L_g$ is the mean loss over sensor $g$'s $n_g$ tiles, then $n_g L_g$ is the \emph{sum} over those tiles, and
\[
  \frac{\sum_g n_g L_g}{\sum_g n_g} = \frac{\text{sum of every tile's loss}}{\text{number of tiles}},
\]
the plain mean over all tiles in the batch: every tile counts equally, whichever sensor it is from. With $n = (3, 3, 2)$ and per-sensor losses 0.05, 0.06,
0.08 (illustrative\tbrk\tagA): $(3 \times 0.05 + 3 \times 0.06 + 2 \times 0.08)/8 = (0.15 + 0.18 + 0.16)/8 = 0.49/8 = 0.0613$, whereas an unweighted mean of
the three would give $(0.05 + 0.06 + 0.08)/3 = 0.0633$.

\subsection{Why this fixes two things at once}

\textbf{BatchNorm.} Every micro-batch now holds all three sensors. Training normalises by the mix's statistics; the running statistics are averages of mixes;
evaluation uses those. The offsets of \cref{sec:finetune:seesaw} are no longer removed in training and reintroduced at evaluation: they are present in both.

\textbf{Gradient balance.} Every optimiser step now contains $3 \times 4 = 12$ EnMAP, 12 PRISMA and $2 \times 4 = 8$ AVIRIS-NG tiles.

\begin{worked}[label={we:finetune:toy}]{Alternating sensors never settle; mixed batches do}
One weight $w$. Sensor A wants $w = 1$ (loss $(w - 1)^2$, gradient $2(w - 1)$); sensor B wants $w = -1$ (loss $(w + 1)^2$, gradient $2(w + 1)$). Learning rate 0.25.

\textbf{Alternating}, one sensor per step, from $w = 0$. Write each step as a map first:
\begin{align*}
  \text{A step:}\quad w &\leftarrow w - 0.25 \times 2(w - 1) = w - 0.5w + 0.5 = 0.5w + 0.5\\
  \text{B step:}\quad w &\leftarrow w - 0.25 \times 2(w + 1) = w - 0.5w - 0.5 = 0.5w - 0.5.
\end{align*}
Then run it:
\begin{align*}
  \text{A:}&\quad 0.5 \times 0 + 0.5 = 0.5\\
  \text{B:}&\quad 0.5 \times 0.5 - 0.5 = -0.25\\
  \text{A:}&\quad 0.5 \times (-0.25) + 0.5 = 0.375\\
  \text{B:}&\quad 0.5 \times 0.375 - 0.5 = -0.3125.
\end{align*}
Where does it end? Compose the two maps: an A step then a B step takes $w$ to $0.5(0.5w + 0.5) - 0.5 = 0.25w + 0.25 - 0.5 = 0.25w - 0.25$. A value that
this pair leaves unchanged satisfies $w = 0.25w - 0.25$, so $0.75w = -0.25$ and $w = -1/3$ after each B step; after each A step it is
$0.5 \times (-1/3) + 0.5 = 1/3$. The weight swings for ever between $+1/3$ and $-1/3$, and where you read it depends on which sensor went last.

\textbf{Mixed}, both losses averaged in every step. The gradient of $\tfrac{1}{2}[(w-1)^2 + (w+1)^2]$ is $\tfrac{1}{2}[2(w - 1) + 2(w + 1)] = 2w$, so
\[
  w \leftarrow w - 0.25 \times 2w = 0.5w.
\]
From $w = 0.5$: $0.25$, $0.125$, $0.0625$, \dots $\to 0$. It settles at the compromise, 0, and stays there.
\end{worked}
Real training is not this extreme (FT02's accumulation already mixed most steps), so the BatchNorm effect is expected to be the larger of the two.

\subsection{What else changed when the next run started}

FT02b was ``FT02's config plus \code{mixed_batches: true}'', from FT02's epoch-8 weights. But \code{init_from} loads \emph{weights only}, and the epoch
counter restarts at 0. So the optimiser state is fresh, the learning rate returns to 0.0001 on a new cosine, and the SAM ramp restarts. At FT02's epoch 8
(index 7 of a 20-epoch cosine):
\begin{align*}
  \eta &= 0.0001 \times \frac{1 + \cos(7\pi/20)}{2} = 0.0001 \times \frac{1 + 0.454}{2} = 0.0001 \times 0.727 = 0.0000727,\\
  \lambda &= 0.5 \times \min(1, 7/10) = 0.35,
\end{align*}
against FT02b's first epoch at $\eta = 0.0001$ and $\lambda = 0.5 \times 0/10 = 0$. If FT02b moves, mixed batches are not the only change.

\begin{sourcenote}[: warm-up is never read]
FT02's and FT02b's configs set \code{lr_schedule.warmup_epochs: 1} and their READMEs say ``1 warm-up epoch''. \code{FoundationTrainer._build_scheduler}
(\file{app/abstract_classes/foundation_trainer.py}, line 517) builds a plain \code{CosineAnnealingLR} stepped once per epoch, and no code reads
\code{warmup_epochs}. There is no warm-up; the first epoch runs at the full rate. \file{docs/lld/finetune.md} now says so.
\end{sourcenote}

\begin{drift}
FT02b (2026-10-16, 2.2 hours, \code{best.pt} epoch 9): validation SAM 3.99 / 4.37 / 5.57$^\circ$ (EnMAP / PRISMA / AVIRIS-NG) against FT02's 4.17 / 5.93 /
6.07; most of PRISMA's drop arrived in the first epoch (4.70). The see-saw was gone by the README's account; the bench gate was missed again, 0.856
against 0.861.\tbrk\tagM{research/experiments/indradhanu-finetune/FT02b-mixed/README.md} \Cref{ch:ftruns} reads the run.
\end{drift}

\section{Joint training: reconstruction and segmentation together}\label{sec:finetune:joint}

\subsection{The loss}

Joint mode adds one term to the mixed reconstruction loss (\code{_mixed_loss}):
\begin{lstlisting}
loss = total / kept                      # reconstruction: L1 + SAM weight x SAM, per tile
if cfg.mode == "joint":
    seg = self._mixed_seg_loss(groups, model)
    if seg is not None:
        loss = loss + cfg.seg_weight * seg   # seg_weight x masked BCE + Dice
\end{lstlisting}
In symbols, with $w_s$ the segmentation weight:
\begin{equation}\label{eq:finetune:joint}
  \mathcal{L} = \underbrace{\frac{\sum_g n_g\bigl(L1_g + \lambda\,\SAM_g\bigr)}{\sum_g n_g}}_{\text{reconstruction, all tiles}}
              \;+\; w_s\,\underbrace{\frac{1}{\abs{\mathcal{G}_\ell}}\sum_{g \in \mathcal{G}_\ell}\mathcal{L}^{\text{seg}}_g}_{\text{labelled groups only}},
\end{equation}
where $\mathcal{G}_\ell$ is the set of sensor groups that carry labels and $\mathcal{L}^{\text{seg}}_g$ is \code{masked_bce_dice} (\cref{ch:seghead}) on group
$g$. In FT03 only EnMAP carries labels, so the second term is EnMAP's segmentation loss. The code's \code{_mixed_seg_loss} returns
\code{sum(losses) / len(losses)} over the labelled groups.

\subsection{Choice 1: the segmentation pass sees the whole mixed batch}

Segmentation needs a second forward pass: the reconstruction pass hides 65\% of the tokens, and a mask predictor must see the whole tile. That second pass
is a danger, because in training mode BatchNorm normalises with \emph{and records} the statistics of every forward call. Suppose the segmentation pass held
only the 3 EnMAP tiles. Every step would then run one mixed pass and one EnMAP-only pass.

\begin{worked}[label={we:finetune:fixedpoint}]{Where the running mean settles with an EnMAP-only second pass}
Mixed batches have mean $a = 0.2$, EnMAP-only batches $b = 0.5$ (illustrative\tbrk\tagA); the running mean starts at 0.2.
\begin{align*}
  \text{step 1, mixed:}&\quad 0.9 \times 0.200 + 0.1 \times 0.2 = 0.180 + 0.020 = 0.200\\
  \text{step 1, EnMAP-only:}&\quad 0.9 \times 0.200 + 0.1 \times 0.5 = 0.180 + 0.050 = 0.230\\
  \text{step 2, mixed:}&\quad 0.9 \times 0.230 + 0.1 \times 0.2 = 0.207 + 0.020 = 0.227\\
  \text{step 2, EnMAP-only:}&\quad 0.9 \times 0.227 + 0.1 \times 0.5 = 0.2043 + 0.050 = 0.2543.
\end{align*}
It keeps climbing. Call the settled value (right after an EnMAP-only pass) $r^\ast$; one full step must leave it unchanged:
\begin{align*}
  \text{after the mixed pass:}&\quad 0.9\,r^\ast + 0.1\,a\\
  \text{after the EnMAP-only pass:}&\quad 0.9\,(0.9\,r^\ast + 0.1\,a) + 0.1\,b = 0.81\,r^\ast + 0.09\,a + 0.1\,b.
\intertext{Set this equal to $r^\ast$ and solve:}
  r^\ast &= 0.81\,r^\ast + 0.09\,a + 0.1\,b\\
  0.19\,r^\ast &= 0.09\,a + 0.1\,b\\
  r^\ast &= \frac{0.09\,a + 0.1\,b}{0.19} = \frac{0.09 \times 0.2 + 0.1 \times 0.5}{0.19}\\
         &= \frac{0.018 + 0.05}{0.19} = \frac{0.068}{0.19} = 0.358.
\end{align*}
The weights $0.09/0.19 = 0.47$ on the mix and $0.1/0.19 = 0.53$ on EnMAP alone show that the stored mean would settle about halfway between them, matching
none of the sensors: the fault FT02b had removed one run earlier.
\end{worked}

So the segmentation pass runs over all eight tiles, through \code{segment_mixed}:
\begin{lstlisting}
def segment_mixed(self, groups: list) -> list:
    # groups: [(x, wavelengths_nm, band_valid), ...] - one per sensor
    z = torch.cat([self._stem(x, wl, bv) for x, wl, bv in groups])   # (8, 32, H, W)
    logits = self.seg_head(self.encoder(self.norm(z)))   # (8, 5, H, W), no keep mask
    # split back by sensor: (3, 5, H, W), (3, 5, H, W), (2, 5, H, W)
    return list(torch.split(logits, [len(x) for x, _, _ in groups]))
\end{lstlisting}
It is \code{forward_mixed} without a keep mask and with the segmentation head in place of the decoder. The PRISMA and AVIRIS-NG rows get logits that no loss
reads; they are in the pass so that BatchNorm normalises, and records, the same mix as the reconstruction pass. With both passes at mean $a$, the fixed point
is $r^\ast = (0.09a + 0.1a)/0.19 = a$: the stored mean is the mix's.

\coursefig[0.95\linewidth]{iw-55-joint-passes}{One mixed batch of 8 tiles goes through the shared BatchNorm and encoder twice: masked, to the decoder, for the
reconstruction loss on all 8; unmasked, to the segmentation head, for the segmentation loss on the 3 EnMAP tiles only (the other 5 rows are normalised but
not scored). Total loss $=$ reconstruction $+ 0.1 \times$ segmentation.}{fig:finetune:joint}

\paragraph{The argument-order trap.} \code{_seg_inputs} returns \code{(pixels, band_valid, wavelengths, targets, known)}; \code{segment_mixed} wants
\code{(x, wavelengths, band_valid)}. The call reorders them:
\begin{lstlisting}
logits = model.segment_mixed([(p[0], p[2], p[1]) for p in parts])   # (x, wl, band_valid)
\end{lstlisting}
Pass \code{p[:3]} instead and the stem receives a \shape{(B, N, H, W)} validity mask where it expects an \shape{(N,)} list of wavelengths. Two neighbouring
functions that order the same three arguments differently is a trap; a short smoke run on real batches is cheaper than a failed three-hour run.

\begin{sourcenote}[: the validator's stale name]
\code{SpectralStemFinetuneConfig._mixed_batches_in_adapt_only} (\file{app/models/training/training_config.py}) now rejects only \code{probe}, with the
message ``mixed\_batches is implemented for modes 'adapt' and 'joint''' (line 231); its name was not changed when joint mode gained \code{segment_mixed}.
The course (IW55) adds that \file{docs/lld/finetune.md} ``still says mode adapt only in four places''; the LLD has since been updated (``mixed\_batches:
adapt or joint''). The source wins, and today the LLD agrees with it.
\end{sourcenote}

\subsection{Choice 2: \texttt{seg\_weight} 0.1, not 1.0}

FT03's README: ``laptop smoke: segmentation loss \textasciitilde 0.4 against reconstruction \textasciitilde 0.05 (L1 $+$ 0.5 SAM)''.\tbrk\tagM{research/experiments/indradhanu-finetune/FT03-joint/README.md}
Where does 0.05 come from? Per-sensor validation in FT02b had L1 about 0.016 and SAM about 4.2$^\circ$. \code{SAMLoss} returns radians:
\begin{align*}
  4.2^\circ &= 4.2 \times 3.1416 / 180 = 0.0733 \text{ rad}\\
  \text{reconstruction} &= 0.016 + 0.5 \times 0.0733 = 0.016 + 0.0367 = 0.0527.
\end{align*}
Now the combined loss \eqref{eq:finetune:joint} at two weights, with segmentation 0.4:
\begin{align*}
  w_s = 1.0:&\quad 0.05 + 1.0 \times 0.4 = 0.45; \quad\text{segmentation share } 0.40 / 0.45 = 89\%\\
  w_s = 0.1:&\quad 0.05 + 0.1 \times 0.4 = 0.09; \quad\text{segmentation share } 0.04 / 0.09 = 44\%.
\end{align*}
At 1.0 the shared encoder would be trained mostly to segment EnMAP, and reconstruction, which is what detection reads, would become a side task. At 0.1 the two
are comparable, as the proposal's ``MAE dominant'' intends. A share of the loss \emph{value} is a rough guide, not the share of the gradient (two terms of
equal size can have gradients of very different size), but an order-of-magnitude imbalance is what it catches.

\subsection{Choice 3: the probe's head as a warm start}

\index{warm start}\code{seg_head_from} names a checkpoint whose \code{seg_head.*} weights replace the fresh head after \code{init_from} has loaded the backbone. FT03 took FT01's
probe head. The plan had said ``the probe is a measurement, not a warm start''; it became both. The FT03 README is careful about what that buys: FT01 ``learned
on F001 features and the encoder has moved since, so it is a warm start, not a fit''.

\begin{drift}
FT03 (2026-10-16, 3.0 hours on the L40S, 10 epochs, \code{sam_ramp_epochs: 0} so the restarted epoch counter does not train epoch 1 on L1 alone): EnMAP test IoU
water 0.841 and cloud 0.815 at \code{best.pt} (epoch 10), above FT01's 0.82 / 0.80; per-sensor SAM 3.92 / 4.21 / 5.48$^\circ$, within 0.1$^\circ$ of its
predecessor; detection unchanged within noise. All three criteria passed.\tbrk\tagM{research/experiments/indradhanu-finetune/FT03-joint/README.md}
\Cref{ch:ftruns} reads the run in full.
\end{drift}

\paragraph{Why joint training did not lift detection, and was not expected to.} The detection score is residual Mahalanobis on the \emph{reconstruction}. The
segmentation head is a separate output; nothing in the score reads it. Joint training could only change detection by changing the encoder, and at
$w_s = 0.1$ the encoder kept reconstructing as before.

\begin{sourcenote}[: no gradual unfreezing]
The proposal's first draft said stage 3 trains with the ``encoder unfrozen from the top stage down'' (gradual unfreezing). \file{finetune_trainer.py} has no
such mechanism: \code{joint} trains every parameter from the first step. The source won, and the proposal was corrected (``gradual top-down unfreezing is not
implemented'').
\end{sourcenote}

\section{How segmentation was meant to help detection}\label{sec:finetune:classaware}

Stage 5 is where the head pays back. Its masks can replace \code{scene_segmentation}'s thresholds, removing cloud edges and shadows before scoring. A second
idea was \emph{class-aware} residual Mahalanobis: give each surface class its own normal. The mathematics behind it is the \emph{law of total variance}\index{law of total variance}.

Let a residual $X$ (one whitened direction) come from land with probability $\pi_1$ and from sea with probability $\pi_2$, with class means $\mu_c$ and
variances $\sigma_c^2$. The pooled mean is $\mu = \sum_c \pi_c\mu_c$. The pooled variance is the mean squared distance from $\mu$; split each class's
distances around its own mean:
\begin{align*}
  \operatorname{Var}(X) &= \sum_c \pi_c\,\E\bigl[(X - \mu)^2 \mid c\bigr]
  = \sum_c \pi_c\,\E\bigl[\bigl((X - \mu_c) + (\mu_c - \mu)\bigr)^2 \mid c\bigr].
\intertext{Expand the square. The cross term $2(\mu_c - \mu)\E[X - \mu_c \mid c]$ is zero, because a class's values average to its own mean:}
  \operatorname{Var}(X) &= \underbrace{\sum_c \pi_c\,\sigma_c^2}_{\text{within classes}} + \underbrace{\sum_c \pi_c\,(\mu_c - \mu)^2}_{\text{between classes}}.
\end{align*}

\begin{worked}[label={we:finetune:totalvar}]{Pooled against per-class scoring on a coastal scene}
A set1-like scene, 75\% sea (illustrative\tbrk\tagA). Land residuals centre on 0 with variance 1; sea residuals centre on 2 with variance 0.25.
\begin{align*}
  \mu &= 0.25 \times 0 + 0.75 \times 2 = 1.5\\
  \text{within} &= 0.25 \times 1 + 0.75 \times 0.25 = 0.25 + 0.1875 = 0.4375\\
  \text{between} &= 0.25 \times (0 - 1.5)^2 + 0.75 \times (2 - 1.5)^2\\
                 &= 0.25 \times 2.25 + 0.75 \times 0.25 = 0.5625 + 0.1875 = 0.75\\
  \text{total} &= 0.4375 + 0.75 = 1.1875.
\end{align*}
Two pixels with residual 3.0: a land target, and an ordinary sea pixel two spreads above the sea's mean. Score each by squared distance over variance:
\begin{align*}
  \text{pooled, either pixel:}&\quad (3.0 - 1.5)^2 / 1.1875 = 2.25 / 1.1875 = 1.89\\
  \text{per class, land target:}&\quad (3.0 - 0)^2 / 1 = 9.0\\
  \text{per class, sea pixel:}&\quad (3.0 - 2)^2 / 0.25 = 1 / 0.25 = 4.0.
\end{align*}
Pooled, they are indistinguishable; per class, the target scores more than twice the sea pixel.
\end{worked}

\begin{drift}
Class-aware scoring was tested and not shipped: water's tight statistics made water noise dominate (set1 water share of false alarms 14\% $\to$ 92\%). Exclusion of
cloud by the head's masks was chosen instead (water was kept, since vessels are in scope), lifting FT03's real-bench pAUC from 0.158 to 0.166 (2026-10-17/18,
\cref{ch:ftruns,ch:shipping}).\tbrk\tagM{docs/proposals/fine-tuning-and-segmentation.md}
\end{drift}

\begin{confusions}
\confusion{Does \texttt{requires\_grad = False} freeze a layer?}{It freezes its \emph{parameters}. BatchNorm's running statistics are buffers and keep updating in
training mode; only \texttt{eval()} stops them.}
\confusion{Evaluation mode turns BatchNorm off?}{It turns its \emph{updating} off and switches it to the stored statistics. It still normalises.}
\confusion{BatchNorm is in every SegFormer block?}{Not in this model. The blocks use LayerNorm. The one BatchNorm is \texttt{self.norm}, in front of the encoder.}
\confusion{Accumulating 4 micro-batches of 8 means BatchNorm sees 32?}{It sees 8. Accumulation adds gradients after four separate forward passes.}
\confusion{Mixed batches mean the stem reads three band lists at once?}{It is still called once per band list. Mixing starts after the stem.}
\confusion{Should validation be mixed too?}{It does not need to be. Validation runs in evaluation mode, so BatchNorm uses its stored statistics, not the batch's; a
sensor-pure validation batch is normalised exactly as a mixed one would be.}
\confusion{Is the normalisation curve updated for the new sensors?}{No, never. The record of trained wavelengths grows; the curve that normalises inputs does not.}
\confusion{Joint means PRISMA and AVIRIS-NG are trained to segment too?}{In FT03 they pass through the segmentation pass but no loss reads their logits; they are there
for BatchNorm. (From FT04 they are taught from confident pseudo-labels, \cref{ch:ftruns}.)}
\confusion{The see-saw is proven to be BatchNorm?}{It is consistent with it. Gradient competition would also produce a see-saw; mixed batches address both, so one run
cannot tell them apart.}
\end{confusions}

\begin{checkyourself}
\item\label{ex:finetune:1} A BatchNorm channel's running mean is 0.10 and training-mode batches all have mean 0.30. Compute it after each of 3 batches, then check
  the third value with the closed form $r_n = 0.9^n r_0 + (1 - 0.9^n)\,x$.
\item\label{ex:finetune:2} What fraction of the stored mean comes from the last 5 micro-batches, and how many micro-batches carry 95\% of it? ($\ln 0.05 = -2.996$,
  $\ln 0.9 = -0.1054$.)
\item\label{ex:finetune:3} The record holds 1000 and 1010\,nm. Extend it with 999.8, 1004.6, 1005.0, 1030.0. What is the record, what does the call return, and is a
  band at 1018\,nm then accepted?
\item\label{ex:finetune:4} With \texttt{batch\_size} 4 and shares 0.40 / 0.35 / 0.25, what are the sub-batch sizes and the effective shares? In the toy problem of
  \cref{we:finetune:toy}, start at $w = 1$ and alternate B, A, B.
\item\label{ex:finetune:5} Redo the running-mean fixed point of \cref{we:finetune:fixedpoint} with mixed mean 0.1 and EnMAP-only mean 0.6. What is it with
  \texttt{segment\_mixed}?
\item\label{ex:finetune:6} A sensor's validation L1 is 0.015 and its SAM 4.5$^\circ$. Compute the reconstruction loss at SAM weight 0.5. With segmentation loss 0.3
  and reconstruction 0.06, what share of the loss is segmentation at $w_s = 1.0$, 0.1 and 0.2?
\end{checkyourself}

\nextchapter{ch:inference}
\endgroup
